\documentclass[a4paper,11pt]{article}
\usepackage{amsmath,amssymb,amsthm}
\usepackage{color,xcolor}
\usepackage{fontspec}
\usepackage[lite,subscriptcorrection,slantedGreek]{mtpro2}
% \DeclareMathSizes{10}{10}{5}{5}
\linespread{1.5}
\usepackage{hyperref}
\hypersetup{
    final,                   % Include links even if document is in 'draft' mode.
    hidelinks,               % Prevent boxes from being drawn around links in some viewers.
    breaklinks=true,         % Allow line breaks in the middle of links
    colorlinks=true,
    linkcolor=black,         % TOC, links to labeled equations and environments
    urlcolor=black!30!blue,  % URLS including links in references  
    anchorcolor=blue,        % I'm not sure what this does.
    citecolor=black!30!blue, % In-line citations  
}
\usepackage[
    noabbrev,               % E.g., "Figure" instead of "Fig."
    capitalise,             % E.g., "Figure" instead of "figure"
    nameinlink              % Make "Figure 1" linked instead of just "1".
]{cleveref}
\setmainfont{Times New Roman}
\newtheoremstyle{sltheorem}
                {}          % Space above
                {}          % Space below
                {\upshape}  % Body font
                {}          % Indent amount
                {\bfseries} % Head font
                {:}         % Punctuation after head
                { }         % Space after theorem head
                {}          % Theorem head spec
% Enable the new "sltheorem" theorem style.
\theoremstyle{sltheorem}
\makeatletter
\def\@endtheorem{\endtrivlist}
\makeatother
\newtheorem{definition}{Definition}
\newtheorem{conjecture}{Conjecture}
\newtheorem{theorem}{Theorem}
\newtheorem{lemma}[theorem]{Lemma}
\newtheorem{corollary}{Corollary}[theorem]
\newtheorem{proposition}{Proposition}
\newtheorem*{remark}{Remark}
\everymath{\displaystyle}
\newcommand{\CollatzReduce}[2]{\Delta({#1};{#2})}
\newcommand\quotient[2]{{{\displaystyle #1}}\big\backslash{{\displaystyle #2}}}
\newcommand\ModSupportSet[2]{{{\displaystyle #1}}-{{\displaystyle #2}}}
\newcommand{\floor}[1]{\left\lfloor#1\right\rfloor}

\begin{document}
\begin{definition}[Collatz Series and Collatz Accumulation]
	Given three arguments $\alpha, \beta, \gamma \in \mathbb{N}$. For number $a \in \ModSupportSet{\mathbb{N}}{\beta\mathbb{N}}$, define a series $\left\{a_{_k}\right\}_{k = 0}^{\infty} \subset \ModSupportSet{\mathbb{N}}{\beta\mathbb{N}}$ and $\CollatzReduce{a}{k}$ with
	\[\begin{aligned}
        a_0 &=a, \\
        \CollatzReduce{a}{0} &=0, \\
        a_{_{k + 1}} \cdot \beta^{\CollatzReduce{a}{k + 1}} &= \alpha\cdot a_{_k} + \gamma,
    \end{aligned}\]
	then the series is called Collatz series or Collatz sequence of $a$, the $\CollatzReduce{a}{k}$ is called Collatz reduce of $a$ at $k$.
    Denote 
	\[\mathcal{A}_n^{m} = \sum_{i=m}^{n}\CollatzReduce{a}{i}\]
	called Collatz accumulation. Specially when $m = 0$, we omit the superscript, that $\mathcal{A}_n = \mathcal{A}_n^{0}$.\newline
	Besides, denote $\mathfrak{A}_n^{m}$ called Collatz complement of $a$ on $k$ with
	\[ \mathfrak{A}_n^{m} = \sum_{i = m}^{n - m - 1}\beta^{\mathcal{A}_i^{m}}\cdot\alpha^{n-m-i-1}. \]
	Also I will omit the superscript $m$ when $m = 0$.
\end{definition}
\begin{proposition}[Consistency]\label{proposition:consistency}
With the same arguments $(\alpha, \beta, \gamma)$, if there are two numbers $a, b \in \ModSupportSet{\mathbb{N}}{\beta\mathbb{N}}$, that $a = b$, then $\forall m > 0$, 
there is $a_m = b_m$ and $\mathcal{A}_m = \mathcal{B}_m$.
\end{proposition}
\begin{proof} Trivial. \end{proof}
\begin{theorem}[Collatz Series's Inductive Formula]\label{theorem:inductive-formula}
Given an odd number $a$, we have
\[ a_k \cdot \beta^{\mathcal{A}_k} = a \cdot\alpha^k + \gamma\cdot\mathfrak{A}_k. \]
\end{theorem}
\begin{proof} TODO \end{proof}
\begin{proposition}[Co-consistency]\label{proposition:co-consistency}
    With the same arguments $(\alpha, \beta, \gamma)$, if there are two numbers $a, b \in \ModSupportSet{\mathbb{N}}{\beta\mathbb{N}}$ that $\forall m \le n$, $\mathcal{A}_m = \mathcal{B}_m$, and $\gcd(\alpha, \beta) = 1$, then $\beta^{\mathcal{A}_n} \mid (a - b)$ and $\alpha^{n} \mid (a_n - b_n)$.
\end{proposition}
\begin{proof} Trivial. \end{proof}
\begin{remark}
for all $m,n \ge 0$, 
\[ a_{m+n} \cdot \beta^{\mathcal{A}_{m+n}^{m}} = a_m\cdot\alpha^n + \gamma\cdot\mathfrak{A}_{m + n}^{m}. \]
\end{remark}
\begin{remark}
    If $\exists n > m \ge 0$ that $a_n = a_m$, 
    then $\exists d \mid (n - m)$ that $a_{m + k \cdot d} = a_{m}$ holds for all $k > 0$.
\end{remark}
\begin{proof} TODO \end{proof}
\begin{remark}
    If $\exists t > 0$ that $a_{t} = a_0 = a$, 
    then $\forall k > 0$, and $s = 1, 2, 3, \cdots, d$, there are
    \[a_{k\cdot t+s} = a_s,\quad \mathcal{A}_{k\cdot t + s} = \mathcal{A}_{s}.\]
\end{remark}
\begin{proof} TODO \end{proof}
\begin{remark}[Multistability]\label{remark:multistability}
	If there exist a subseries that $\lim_{n\to\infty} \frac{\mathcal{A}_n}{n} = \mu$, let $\eta = \frac{\beta^\mu}{\alpha}$, then for the subseries, we have
    \[\mathfrak{A}_n = \left\{\begin{aligned}
    & \alpha^{n - 1} \cdot \frac{1 - \eta^n}{1 - \eta}, &\eta \ne 1, \\
    & \alpha^{n-1} \cdot n, &\eta = 1,
    \end{aligned}\right.\]
    and then
	\[a_n = \left\{\begin{aligned}
	& \frac{a}{\eta^n} + \frac{\gamma}{\alpha} \cdot \frac{\eta^n - 1}{1 - \eta}, &\eta \ne 1, \\
	& a + \frac{\gamma}{\alpha}\cdot n, &\eta = 1.
	\end{aligned}\right.\]
\end{remark}
\newpage
From the recursion fomula
\[ a_{k+1} \cdot \beta^{\CollatzReduce{a}{k + 1}} = \alpha \cdot a_{k} + \gamma \]
we have
\[ a_k \equiv - \gamma \cdot \alpha^{-1} \pmod{\beta^{\CollatzReduce{a}{k + 1}}}. \]
This means
\[ a_k \cdot \beta^{\mathcal{A}_{k}} \equiv - \gamma \cdot \alpha^{-1} \cdot \beta^{\mathcal{A}_{k}} \pmod{\beta^{\mathcal{A}_{k+1}}}, \]
means
\[a \cdot\alpha^k + \gamma\cdot\mathfrak{A}_k \equiv - \gamma \cdot \alpha^{-1} \cdot \beta^{\mathcal{A}_{k}} \pmod{\beta^{\mathcal{A}_{k+1}}}, \]
means
\[a \cdot\alpha^k \equiv - \gamma \cdot \alpha^{-1} \cdot \beta^{\mathcal{A}_{k}} - \gamma\cdot\mathfrak{A}_k \pmod{\beta^{\mathcal{A}_{k+1}}}. \]
Then we have
\[ a \equiv - \gamma \cdot \alpha^{-(k + 1)} \cdot (\alpha\cdot\mathfrak{A}_{k} + \beta^{\mathcal{A}_{k}}) \pmod{\beta^{\mathcal{A}_{k+1}}}. \]
If there is an subseries saticifies the condition of \textbf{Multistability}, i.e. $\lim_{n\to\infty} \frac{\mathcal{A}_n}{n} = \mu$ and $\eta = \frac{\beta^\mu}{\alpha}$, then we have
\[a \equiv \left\{\begin{aligned}
& - \gamma \cdot \alpha^{-1} \cdot \frac{1 - \eta^{n}}{1 - \eta}, &\eta \ne 1 \\
& - \gamma \cdot \alpha^{-1} \cdot n , &\eta = 1
\end{aligned}\right. \pmod{\beta^{\mathcal{A}_{n}}}.\]
Here, negative exponents and fractions should be understood in the sense of modulus (integer field).
\end{document}